% 本文件由 LaTeX 排版 Agent 根据有效 translation.json 自动生成。
% author=Augustine of Hippo; work_id=1316; title=论对亡者的照顾

\chapter{论对亡者的照顾}

% structure: p0001 source=h1 target=omitted reason=page-title-already-rendered-by-parent

摘自\WorkTitle{订正录}卷二，第64章。\par

\WorkTitle{论对亡者应尽的照顾}一书，系因有人写信问我：若将亡者的遗体埋葬于某位圣人的纪念地，是否对其死后有益，遂撰此书。书之开篇如下：\newline{}\Quote{久仰尊驾圣德，我可敬的同道主教宝利诺。}\par

1. 我可敬的同道主教宝利诺，久欠答复尊驾之函；自您经由我们最虔诚的女儿福乐拉家中之人来信，问我：人死后，将其遗体葬于某位圣人的纪念地，是否对其有益。原来那位寡妇曾向您请求，为她死于贵地的儿子办理此事，而您回信安慰她，并告知她有关那位忠信青年基乃吉的遗体之事：她以慈母般的虔诚心愿所求之事，业已照办，即将其安放在最蒙福的精修圣人斐理斯的圣堂内。为此，您亦藉同一位信差致信于我，提出同样的问题，并求我答复对此事的看法，同时您也不隐瞒自己的意见。因为您说，依您之见，那些为去世亲友操办此事的宗教及信德之心，并非空洞无谓的举动。您更补充说，整个教会惯于为亡者祈祷，这不可能毫无功效：因此可以进一步推断，若因亲友的信德，为埋葬其遗体而选择了这样的场所，使得圣人们也能以这种方式寻求的帮助得以彰显，这对人死后是有益的。\par

2. 但既然如此，如何使宗徒所说的\Quote{我们众人都要站在基督的审判台前，好使每人按他在身体上所行的，或善或恶，领取相当的报应}与这一意见不相矛盾？您表示对此不甚明了。因为这句宗徒之言告诫的是应在死前做那些死后有益之事，而非在此时才去领受一个人在死前所行之事。然而，此问题可如此解决：存在某种生活方式，在人活于肉身之时即已获得，使得这些事物可能对亡者有所帮助；因此，正是\Quote{按身体所行的}，他们才得以借那些在离世后为他们虔诚举行的事而获得助益。因为有些人，这些事对他们毫无帮助——即当这些事是以如此恶劣功绩之人为对象时，他们不配得到这类帮助；或当这些事是以如此优良功绩之人为对象时，他们并不需要这种帮助。因此，每个人在肉体中所过的生活方式，决定了在他离世后那些为他虔诚举行的事，是否对他有益或无效。因为若在今生未获得使这些事物有益的功德，死后寻求便属徒劳。由此，一方面，教会或亲友的关怀为亡者所能提供的任何宗教行为并非毫无意义；另一方面，每个人仍\Quote{按他在身体上所行的，或善或恶}领受报应，主照各人的行为予以回报。因为，为使这些施予之物能在死后对他有益，那是在他活于肉体中的一生中就已获取的。\par

3. 或许您对此疑问已因我这简短的答复而满足。但请稍待，听我说说另有何种想法使我动心，我以为有必要对之答复。在玛加伯书中，我们读到为亡者献祭的记载。然而，即使古经中从未读到，整个教会在这一习俗中明白显示的权威也不容小觑，即在司祭于天主的祭台上向主天主献上的祈祷中，也包含了对亡者的托付。但关于亡者灵魂是否因遗体安放之处而获得某种益处，则需更仔细地探讨。首先，若尸体不得埋葬，是否会造成人死后灵魂受苦或加深痛苦，这不应依于普遍看法，而应依据我们宗教的圣经。因为我们不应相信，如维吉尔的诗句所言，未埋葬者不得渡过冥河：因为据说\par

无人能渡阴森之岸与汹涌之水，\newline{}除非其骨骸已在合适之所安息。\newline{}\par

谁能以基督徒之心相信这些诗意的虚构？主耶稣为使基督徒在敌人手下——敌人可控制其身体——能无所畏惧地安息，曾断言他们连一根头发也不会损坏，并劝勉他们不要害怕那些杀害身体后便不能再做什么的人。关于此点，我在\WorkTitle{天主之城}第一卷中，我想已充分论述，以挫败那些人的锐气——他们将罗马近来遭受的野蛮蹂躏归咎于基督时代，还以此责难我们，说基督未来援助祂自己的人。当我们回答说，信友的灵魂已按他们信德的功能受到祂的保护时，他们便以遗体未得埋葬来嘲弄我们。因此，关于埋葬的一切问题，我以这样的言辞作了阐述。\par

\Quote{但是}（我说）在这样的死人堆中，难道他们甚至不能被埋葬吗？虔诚的信德也并不过于畏惧这一点，因为它坚信所预言的：甚至连吞噬的野兽也不能妨碍身体的复活，因为头上一根头发也不会丧失。真理绝不会说：\BibleQuote{不要怕那些杀害身体，却不能杀害灵魂的；}如果仇敌无论对被杀者的身体做什么都能妨碍来世的话。除非有人如此荒唐，竟主张在死前不应怕他们，以免他们杀害身体，但在死后却应怕他们，以免他们杀害身体后，不许身体被埋葬。那么基督所说的：\BibleQuote{杀害身体以后，再没有能力做什么的，}难道是假的吗？如果他们能对死尸做那么大的事？绝无此事，真理所说的话不会是假的。因为所说的是，他们在杀害时做了一些事，因为在身体被杀害时有感觉，但此后他们再不能做什么，因为身体被杀后没有感觉。许多基督徒的身体没有被泥土覆盖，但没有一个身体与天地分离，祂以自己的临在充满天地，祂知道从哪里复活祂所创造的。的确，在圣咏中记载：\BibleQuote{他们把你仆人的尸体，给天空的飞鸟作食物，把你圣人的肉，给地上的野兽作食物；在耶路撒冷四周，他们把他们的血像水一样倒出，没有人埋葬他们。}但这更多是为了突出做这些事的人的残酷，而不是为了表示受这些事的人的不幸。因为在人眼中，这些事可能显得艰难可怕，然而\BibleQuote{在主眼中，祂圣人的死是珍贵的。}所以，所有这些事——丧事的料理、埋葬的安置、葬礼的排场——更多是为了安慰活人，而不是为了帮助死者。如果昂贵的葬礼对恶人有好处，那么简陋或没有葬礼就会损害善人。那位穿紫红袍的富翁在世人眼中得到了他家人非常体面的葬礼；但那位满身疮痍的乞丐却得到了天使的服侍，更为体面；天使没有把他抬进大理石坟墓，而是把他高高举入亚巴郎的怀中。所有这一切都被那些人所嘲笑，我们正是为了这些人而辩护\WorkTitle{天主之城}：然而，就连他们自己的哲学家也鄙视葬礼的关怀；并且常常，整个军队在为自己的世俗国家而战死时，并不在乎他们之后会躺在哪里，或成为什么野兽的食物；诗人们也得以在喝彩中谈论此事\par

虽然他没有骨灰罐安息，\newline{}他的覆盖物是穹苍。\par

他们更不应该夸耀基督徒的身体未被埋葬，因为基督徒的肉身及其所有肢体，不仅从尘土中，而且从其他元素中，甚至从这些消失的尸体所退隐的最隐秘的弯曲处，被重新塑造，并确据将在顷刻之间复原，完全如初，按照祂的应许？\par

5. 然而这并不表示，亡者的遗体——尤其是义人、忠信者的遗体——应当被轻视并丢弃；他们的灵魂曾神圣地使用这些身体，作为成就一切善行的工具和器皿。因为如果父亲的衣裳、戒指及诸如此类之物，对于他所遗留的人更为珍贵（他们对父母的感情越深），那么肉体本身就更不应遭鄙弃；我们与这肉身的结合，远比我们穿戴的任何衣物都更为亲密、更为紧密。因为这些（身体）并非外在装饰或辅助之物，而是属于人的本性自身。因此古时义人的葬礼受到虔敬的照料，丧礼得以举行，坟墓也予以预备；他们尚在人间时，也感人地就自己的埋葬甚至遗骸迁移，嘱咐其子女。\Person{多俾亚}也因埋葬死者而蒙天主恩宠，有天使作证称赞他。主自己也将要第三天复活，祂亲自传扬并令人传扬一位虔诚妇女的善行，她用珍贵的香膏抹在祂身上，为祂的安葬而做；那些从十字架上领受祂的身体，并细心、恭敬地将祂包裹安放在坟墓里的人，也在福音中得着赞美的记载。然而这些权威事例并非要我们以为死者的身体还有知觉；而是要表明，天主的眷顾（祂更喜悦这样的孝爱之事）也亲自眷顾死者的身体，为要藉此坚固我们复活的信德。由此我们也能健康地学到：若对无生命的肢体尽任何义务和殷勤，都不致在天主面前丧失，那么我们对有生命、有知觉的人所行的施舍，将得何等的赏报呢！的确还有别的事，圣祖们在谈论其遗体的安置或迁移时，是愿意人们理解为他们藉着先知之神所说的；但此处不是详论这些事的时候，因我们所说的已足够了。但如果缺乏生活所需之物，如食物和衣服，即使缺乏带来极大的痛苦，也不能摧毁善人承受和忍耐的刚毅勇气，也不能从心中根除孝爱，反而因锻炼而使其更富有成果；那么，缺乏通常用于照料丧葬和安置亡者遗体的物品，又岂能使他们不幸？因为如今他们已在敬虔者隐秘的居所中安息了。因此，在那座大城或其他城镇遭蹂躏时，基督徒的遗体缺乏了这些（丧葬之物），活着的人并无过失，因为他们无力提供；死去的人也并无痛苦，因为他们毫无知觉。以上是我对安葬之根据和理由的看法。我因而从我的另一部书移录于此，因为重述此处比另换一种方式表达同一事更为便捷。\par

6. 若果真如此，那么在圣人的纪念所旁为遗体提供安葬之处，无疑是对朋友遗骸的一份良好人情之标志：因为若埋葬本身是宗教性的，则考虑埋葬在何处也不能不是宗教性的。虽然如此，为了显示生者对所爱者的孝爱之心，这些安慰是可取的；但我看不出它们对死者有何助益，除非是这样：当生者想起他们所爱之人的遗体安放之处时，便藉着祈祷将亡者托付给那些同一位圣人，这些圣人作为主保，已将他们收纳在自己的照管之下，在天主前帮助他们。即使他们不能将遗体安葬在这些地方，仍然能这样做。但那些特别著名的死者坟墓被称为\Quote{纪念所}或\Quote{纪念碑}的唯一理由，是它们唤起记忆，并通过提醒使我们想起那些因死亡而从生者眼前消失的人，以免他们也因遗忘而从人的心中消失。因为\Quote{纪念}一词最清楚地表明了这一点，而\Quote{纪念碑}则得名于\Quote{提醒}，即唤起记忆。因此，希腊人也把我们所称的\Quote{纪念所}或\Quote{纪念碑}称为\OriginalTerm{纪念所}{\GreekText{μνημεῖον}}：因为在他们的话语言中，记忆本身（我们借以记住的能力）被称为\OriginalTerm{记忆}{\GreekText{μνήμη}}。因此，当心灵忆起一位挚友的身体埋葬之处，随即想起一个因殉道圣人之名而变得可敬的地方时，它便以衷心的追忆和祈祷之情，将灵魂托付给那位同一位殉道圣人。当信友中最亲爱的人向亡者表示这种情感时，无疑对死者有益；这些死者身在肉身时，曾配得这些事物在今生之后对他们有益。但即使某种需要由于完全缺乏条件，不允许遗体被埋葬，或埋葬在那些地方，仍不应停止为亡者灵魂的祈祷：这些祈祷应当是普世教会中所有去世的基督徒，即使不提及他们名字，在普遍的纪念之下，教会已亲自担负起来；目的是使那些在这方面缺少父母、子女或任何亲戚朋友的人，也能由那共同慈母（教会）为他们提供同样的祈祷。但如果缺乏这些以正确信德和孝爱为亡者所作的祈祷，我认为无论他们的遗体被安葬在何等神圣的地方，对他们的灵魂也毫无益处。\par

7. 所以，当一位有信德的母亲为她有信德的已故儿子恳求，愿将他的遗体安放在一位殉道者的圣殿中时，因为她相信他的灵魂将因殉道者的功德而得助益，这个信仰本身便是一种祈求，且若有所裨益，便已获益。而她心中常思念同一墓穴，并在祈祷中更加将她儿子交托，这有助于亡者的灵魂，不是靠死者的安葬之地，而是靠对那地方的记忆所生发的、母亲鲜活的爱心。因为立即想到谁被交托、交托给谁，这确实触动祈祷者虔敬的心，且并非无益的情感。同样，在向天主的祈祷中，人们用身体的肢体做出求恳者所行之事：屈膝、伸臂，甚至俯伏于地，以及一切外在可见的动作；虽然他们的不可见的意志和心中的意向天主已洞悉，祂并不需要这些标记来向祂开启任何人的心意；但人借此更能激励自己更谦卑、更热切地祈祷和叹息。我不知为何，虽然这些身体的动作非有先前的心念驱动不能做出，然而当它们外在可见地做出来时，那内在不可见的、使之做出的心念反而增强；因此，那先于动作的心之情感，因动作的完成而增长。但若有人被束缚，甚至被捆绑，以致不能用肢体做这些事，这并不意味着内在的人不祈祷，不在天主眼前、在其最隐秘的内室中、怀着痛悔俯伏于地。同样，当人为亡者的灵魂向天主祈求时，他安放死者的遗体之处确实关系重大，因为先前的情感选择了一个神圣的地点，而在遗体安放后，对那神圣地点的回忆更新并增强了先前的情感；然而，即使他无法将他所爱者埋葬在他虔敬之心所选择的那个神圣之地，他也绝不应在交托亡者时停止必要的祈求。因为无论逝者的肉身安卧何处或不安卧何处，灵魂需要安息并必得安息；因为灵魂在离开肉身时带走了知觉，没有它，人如何存在、处于好的或坏的状态便无关紧要；它并不期望从那肉身获得生命的助益，因为它曾给予那肉身生命，又在离开时收回，并在回归时归还；因为不是肉身给灵魂，而是灵魂给肉身带来甚至复活本身的价值，无论是为了惩罚还是为了荣耀而复活。\par

8. 我们在由欧瑟比以希腊文撰写、鲁菲努斯译为拉丁文的\WorkTitle{教会史}中读到，在高卢，殉道者的尸体被抛给狗，而狗吃剩的残骸和死者的骨头，甚至被火烧尽，直至完全销化；骨灰被撒在罗讷河中，以免留下任何可供纪念的东西。此事必是天主所允许，别无其他目的，为使基督徒在明认基督时，既轻视此生，便更应轻视安葬。因为对殉道者尸体施行的这种野蛮暴行，若丝毫能伤害他们，损害他们最胜利之灵魂的蒙福安息，则天主绝不会允许此事发生。所以，这确实宣告了：主说\BibleQuote{你们不要害怕那杀害肉身，以后不能再做什么，}并非指祂不允许他们向祂的追随者死后的身体做什么；而是指无论他们被允许做什么，都不能减损逝者基督徒的福乐，不能触及他们死后仍活着的知觉；也不能损害身体本身，不能减少它们复活时完整无损的任何部分。\par

9. 然而，由于人心的那种情感——\BibleQuote{没有人会恨自己的身体}，如果人们有理由知道，死后他们的身体将缺少任何在各自民族或国家中习惯的葬礼所要求的东西，他们就会像凡人一样忧愁；而死后并不触及他们的事，他们在死前却为自己的身体恐惧：因此我们在\WorkTitle{列王纪}中看到，天主藉一位先知威胁另一位违背祂话的先知，说他的尸身不得葬入祖坟。圣经如此记载：\BibleQuote{上主这样说：因为你违抗了上主的话，没有遵守上主你的天主吩咐你的命令，竟然返回去，在上主吩咐你不可吃饼、不可喝水的地方吃了饼、喝了水，所以你的尸身不得葬入你祖先的坟墓。}现在，如果按照福音来考虑这个惩罚有何意义，因为我们从福音中得知，身体被杀之后，不必担心无生命的肢体还会遭受什么，那么它甚至不能被称为惩罚。但如果我们考虑人对自身血肉的天然情感，那么他活着时可能会因那死后毫无知觉的事而感到恐惧或悲伤：这是一个惩罚，因为内心因即将发生在自己身体上的事而痛苦，尽管这事发生时他不会有感觉。正是为此，主乐意惩罚祂的仆人，他并非出于自己的悖逆而藐视遵行祂的命令，而是因别人的谎言欺骗，以为自己是在服从，实际上并未服从。因为不应认为他被野兽的牙齿咬死，灵魂就被拽到地狱的酷刑中；况且杀死他的那头狮子还看守着他的尸体，并且他骑的驴也安然无恙，那头凶猛的野兽竟无畏地守在他主人的尸体旁。通过这一奇妙的标记，似乎显示，这位天主的人，毋宁说是在今生被惩戒至死，而非死后受罚。关于这事，宗徒在提到许多人因某些过犯而生病、死亡时说：\BibleQuote{我们若是先省察自己，就不至于受审判了。但我们受审判，正是被主管教，免得我们和世界一同被定罪。}至于那位引诱他的先知，实际上非常尊敬地将他葬在自己的墓中，并安排把自己的骸骨葬在他旁边；希望藉此使自己的骸骨免受侵犯，因为按照那位天主的人的预言，犹大王\Person{约史雅}将在这土地上掘出许多死者的骸骨，并用这些骸骨玷污那些为偶像设立的亵圣祭台。他保存了那位在三百多年前预言这些事之先知的坟墓，并且因他的缘故，那位诱惑他的先知的坟墓也没有被侵犯。也就是说，因为那种使\BibleQuote{没有人会恨自己的身体}的情感，这位用谎言杀害自己灵魂的人，为自己的尸身作了打算。因此，出于每个人对自己血肉的这种自然爱，对于前者，得知自己不能葬在祖坟里是一种惩罚；对于后者，则是一种关切，预先安排好自己的骸骨得以保全，如果他躺在那个坟墓无人敢侵犯的人旁边。\par

10. 为真理而斗争的基督殉道者克服了这种情感；他们既然在活着时感受到的酷刑中都不能被征服，那么他们蔑视死后不再有感觉的事，也就不足为奇了。天主无疑能够（正如祂不让那杀死先知的狮子进一步触碰他的身体，并使那凶手成为守护者那样），我说，祂能够使被扔给狗的祂自己的被杀尸体免遭狗噬；祂能够以无数方式阻止人们的暴怒，使他们不敢焚烧尸身、扬撒骨灰：但是，让这个经验也不缺少于各种试探的多样性是合适的，免得那不愿为保全身体的生命而对迫害的残暴让步的宣认之勇气，还要为坟墓的荣誉而焦虑不安；总之，免得对复活的信德害怕身体的销毁。因此，让这些事情也被允许是合适的，这样，即使在如此可怖的榜样之后，那些在宣认基督的热忱中的殉道者，也成了这一真理的见证，在其中他们学到了：杀害他们身体的人，\BibleQuote{此后不能再做什么了}。因为，无论他们对死尸做什么，终究什么也做不了，因为在那毫无生命的肉身中，那离去的人不可能有任何感觉，而那创造者也不会失去任何一部分。然而，当这些事情正在对被杀者的身体进行时，虽然殉道者并不为此害怕，并勇敢地忍受，但在弟兄们中间却有极大的悲伤，因为他们没有机会向圣人的遗体致以最后的敬意，而且（正如同一历史所记载的）残酷哨兵的看守也不允许他们秘密地取走任何部分。因此，那些被杀的人，在肢体撕裂、骨头焚烧、骨灰抛散中，感觉不到痛苦；但这些没有尸体可埋葬的人，却因怜悯他们而遭受极度痛苦的折磨；因为那些死者毫无感觉的事，这些人在某种程度上为他们感觉到了，而且此后死者不再受苦，这些人在悲哀的同情中为他们受苦。\par

11. 关于我所提到的可悲的同情心，那些为\Person{撒乌耳}和\Person{约纳堂}的枯骨施以埋葬之仁慈的人受到赞美，并蒙\Person{达味}王祝福。但这算是什么仁慈呢？是给予那些毫无感觉之人的吗？或者，这或许应当追溯到一个地狱河流的幻想，认为未被埋葬的人无法渡过去？愿基督徒的信德远离此念！否则，那么多的殉道者遭遇了最坏的情形，因为他们的身体无法被埋葬，而真理欺骗了他们，当它说：\BibleQuote{你们不要害怕那些杀害身体，然后不能再做什么的人，}如果这些人竟能对他们行如此大的恶，以致阻碍他们渡往所渴望的地方，那么真理就欺骗了他们。但是，因为这毫无疑问是极其荒谬的，并且身体的埋葬既丝毫不伤害信徒，也不给不信者带来任何益处；那么，为何那些埋葬\Person{撒乌耳}和他儿子的人被称为行了仁慈，并为此被那位虔诚的君王祝福？正是因为这是一种良好的情感，富有怜悯之心的人的心被触动，当他们为他人死去的身体感到悲伤时，由于那从未有人憎恨自己肉体的情感，他们不愿自己死后自己的身体被如此对待；而他们希望别人在自己毫无感觉时为自己所做的，他们就在自己还有感觉时，为那些现在毫无感觉的人去做。\par

12. 有一些关于显现或异象的故事被讲述，它们似乎为本次讨论引入了一个不应轻视的问题。也就是说，据说死人在某些时候，或是在梦中，或是用其他方式，向那些不知道他们未埋葬的尸体在何处的活人显现，并指出那个地方，告诫他们应当给予所缺的埋葬。如果我们回答这些事情是假的，我们将被认为无耻地反驳某些忠实信仰者的著作，以及那些向我们保证这些事情确实发生在他们身上的人的感觉。但是应当回答，这并不必然导致我们认为死人对此有知觉，因为他们出现在梦中说、指明或请求这些事。因为活人也常常在睡眠中向活人显现，而他们自己并不知道他们显现了；并且梦者被告知，他们梦见了什么，即，在他们的梦中，说话者看见他们做或说某事。那么，如果有可能一个人在梦中看见我向他指示某事已经发生，甚至预告某事将要发生，而我对此完全不知情，并且完全不在乎他梦什么，也不在乎他醒着时我睡着，或他睡着时我醒着，或我们是否同时都醒着或睡着，在他梦见我的那个时刻；那么，死人对此毫无知觉和感觉，却仍被活人在梦中看见，并说一些那些醒来时发现是真的话，这有什么奇怪呢？那么，我认为这是通过天使的行动实现的，无论是从上许可还是命令，使得他们似乎在梦中说一些关于埋葬他们尸体的话，尽管这些尸体所属的人对此完全无知。现在，这有时是有益地完成的；或为了给幸存者某种安慰，那些死者的相似形象在他们的梦中出现；或者通过这些告诫，人类种族可能对埋葬的人道行为有所关注，尽管这对死者没有帮助，但轻视它是有可责的不敬虔。然而有时，通过骗人的异象，人们被抛入大错中，他们应受此苦。例如，如果有人在梦中看到\Person{埃涅阿斯}在诗意的虚假中被说成是在阴间所见的情景：并且有一个未埋葬之人的形象向他显现，说出像\Person{帕利努鲁斯}对他说的话那样的话；当他醒来时，他发现在那个地方有尸体，正如他在梦中听到所说，它未埋葬，并被告诫和请求在他找到时埋葬它；因为他发现这是真的，他就相信埋葬死人是特意为了让他们的灵魂能够到达那些地方，从那里他梦见未被埋葬之人的灵魂被地狱的法律禁止通过：那么，他这样相信这一切时，不是大大偏离了真理之路吗？\par

13. 然而，人性如此软弱：当一个人在梦中看见一个死者，他便以为所见的是灵魂；但当他同样梦见一个活人，他毫不怀疑所显现给他的并非灵魂或身体，而是一个人的肖像；仿佛对于死者，也同样可能在他们不知不觉中，不是他们的灵魂，而是他们的肖像显现在睡者面前。的确，当我们在\Place{米兰}时，我们听说有一人被人凭其已故父亲的借据要求偿还债务，而这笔债务他父亲生前已还清，儿子却不知情。于是此人非常忧愁，惊异于他父亲临终时，在他立遗嘱时，竟未告知他所欠的债。就在他极度焦虑之际，他的父亲在梦中显现给他，告知他何处有解除该借据的收据。这青年找到并出示后，不仅驳斥了假债的不当要求，还取回了父亲还款时未曾收回的借据。此处，人们以为人的灵魂关心其子，并在睡梦中来到他那里，教导他所不知道的事，从而解除他的大烦恼。但大约在同一时期，我们听说在\Place{迦太基}，修辞学家\Person{欧洛吉乌斯}（他曾是我在该艺术上的弟子），据他本人在我们回到\Place{非洲}后告诉我们，正在向弟子们讲授\WorkTitle{西塞罗的修辞学著作}。当他查看次日要讲授的读物部分时，遇到一段文字，不能理解，心中烦恼，几乎不能入睡。那夜，他在梦中，我向他解释了那他不明白的段落；不，不是我，而是我的肖像，而我对这事全然不知，也许远在海外，正在做别的事，或正做别的梦，毫不关心他的烦恼。这些事情如何发生，我不知道；但无论如何发生，我们为何不相信：一个人在梦中看见一个死者，正如他梦见一个活人那样，是同样发生的呢？两者无疑都是，既不知也不关心谁、在何处、何时梦见他们的肖像。\par

14. 同样，有些清醒之人，其感官受到扰乱的异象，例如疯狂的人或任何方式癫狂者，也如同梦境：因为他们也自言自语，仿佛在向真正在场的人说话，同时与不在场者和在场者交谈，这些人的肖像，无论他们是活人或死者，他们都能感知。然而，正如活着的人并不知道自己被他们看见并与之交谈；事实上，他们本人并未真正在场，也未自己发言，只是通过这些感官受扰者，被此类虚幻的异象所影响；同样，那些已离世的人，在如此受影响的人面前显现为在场，而他们本人并不在场，并且无论任何人看见他们的肖像，他们自己全然不知。\par

15. 与此类似的情况还有：当人的感官比睡眠时更深地处于停滞状态时，他们也会被类似的异象所占据。因为在他们面前也会出现活人和死人的形像；但随后，当他们恢复知觉时，他们声称所见到的任何死人，都被认为是确实与他们同在；而听到这些事的人，却不注意同样也看到了某些活人的形像，而他们却是缺席且无意识的。有一个名叫\Person{库尔马}的人，来自\Place{图利姆镇}，靠近\Place{希波}，是市议会中一个贫穷的成员，几乎不能胜任该地双人执政官的职务，而且是个十足的乡下人。他生了病，所有感官都陷入恍惚，几乎像死人一样躺了好几天：只有鼻孔里微弱的呼吸，用手去探才勉强感觉到，几乎表明他还活着，这才使他免于被当作死人埋葬。他一动不动，滴水不进，无论是眼睛还是其他身体感官，都对任何接触到的刺激毫无知觉。然而他看见了许多像梦一样的东西，过了好多天他终于醒来时，他讲述了自己所见。首先，在他睁开眼睛后不久，他说：\Quote{派人到铁匠\Person{库尔马}家里去，看看那里发生了什么事。}有人去了那里，发现铁匠正是在他苏醒、几乎可以说是从死里复活的那个时刻死去的。然后，在围观者急切地倾听时，他告诉他们，那人如何被命令带去（指死者），而他自己则被释放了；他听到在他返回的那个地方有人说，被命令带到那个死人地方的，不是市议会的\Person{库尔马}，而是铁匠\Person{库尔马}。在这些如梦的异象中，在他所见的那些按各自功过被对待的死人中，他也认出了一些他活着时认识的人。我或许会相信他们就是那些人本人，要不是他在这个看似梦境中，也见到了至今仍活着的一些人，即他所在教区的几位神职人员，当地的一位司铎告诉他，他应当由我在\Place{希波}为他领洗，他说这事也确实发生了。因此，他在这个异象中见到了司铎、神职人员、我本人，这些人尚未去世；后来他又见到了死人。为什么我们不能认为他是以同样的方式见到后者的呢？也就是说，两者都是缺席且无意识的，因此不是他们本人，而是他们的肖像，如同地点的肖像一样？他看到了一块土地，那里有那位司铎和神职人员；还有\Place{希波}，他在那里似乎由我领了洗。但在这些地点，当他以为自己在那里时，他实际上并不在那里。因为当时那里发生什么，他并不知道；如果他确实在那里，他肯定会知道。因此，所见到的景象并不是事物本身如何呈现，而是以某种事物的形象投射出来的。总之，在他所见的许多事情之后，他还讲述了如何被带入天堂，以及在那里，当他被释放返回自己的家时，有人对他说：\Quote{去吧，领洗，如果你愿意进入这有福之地。}于是，他受到提醒应由我为他领洗，他说已经领过了。与他谈话的那人回答说：\Quote{去吧，真正地领洗，因为你只是在异象中看到了（领洗）。}此后，他康复了，前往\Place{希波}。复活节临近，他与许多我们素不相识的人一起，作为望教者报了名；他也没有把异象告诉我或我们这边的任何人。他领了洗，圣周结束后返回了自己的地方。过了两年多，我才得知整个事情：起初是通过我的一位也是他的朋友，在我家用餐时我们谈论一些类似事情时得知的；然后我亲自让他本人讲述这个故事，在场有几位他的诚实的同乡作证，关于他那奇异的病，他如何像死人一样躺了许多天，关于那个铁匠\Person{库尔马}，我上面已经提到，以及所有这些事情；他在讲述时，他们回忆起来并向我保证，他们当时也亲耳听他讲过。因此，正如他看到了自己的领洗、我本人、\Place{希波}、圣堂和圣洗池，并非真实本身，而是某种事物的肖像；同样，他也以同样方式看到了某些其他活着的人，而这些活人毫无知觉；那么，那些死人为何不也以同样方式看到，而这些死人毫无知觉呢？\par

16. 为什么我们不应相信这些是凭着天主上智安排而运作的天使的行动？天主利用善与恶，依照祂审判的不可测的深渊，无论是为了教导凡人的心智，或是欺骗它们，或是安慰它们，或是惊吓它们：这一切都按着祂对每人或显示慈悲，或施行惩罚，而教会并非无意义地歌唱\BibleQuote{慈悲与审判}。各人可按自己意愿接受我所说的。如果死者的灵魂参与生者的事务，并且当我们看见他们时，他们在睡眠中对我们说话；且不说别人，就说我自己，我虔诚的母亲每晚必来探望我，那位母亲曾追随我跋山涉水，只为与我同住。切莫认为她因更幸福的生命而变得残忍，以至于当我内心忧愁时，她竟不安慰她以独一之爱所爱的儿子，她从未希望看到他悲伤。但圣咏在我们耳中唱的是真实的：\BibleQuote{因为我的父亲和母亲抛弃了我，但主却收纳了我。}那么，如果我们的父母抛弃了我们，他们又怎能参与我们的忧虑和事务呢？但如果父母不参与，还有哪个死者知道我们的所作所为或所受的苦呢？依撒意亚先知说：\BibleQuote{因为祢是我们的父亲；亚巴郎不认识我们，以色列也不承认我们。}如果如此伟大的圣祖们对他们后裔子民的事一无所知，他们（他们因信天主而获许这子民将从他们的后裔生出）又怎会参与生者的事务和行为，无论是为了知道或帮助？我们怎能说那些在灾祸紧随死亡之前去世的人是有福的，如果死后他们仍然感受到人生灾难中发生的任何事？或者我们这样说是否错了，认为他们安息，而活人的不安生活却使他们焦虑？那么，天主向最虔诚的\Person{约史雅}王所应许的大恩是什么？就是让他先死，免得看见祂警告要降临那地方和那子民的灾祸。天主的这些话如下：\BibleQuote{上主以色列的天主这样说：关于你听到的话，你在我面前畏惧，当你听到我论及这地方和其中居民所说的话，它将荒废并被诅咒；你撕裂了你的衣服，在我面前哭泣，我垂听了你，万军的上主说：不如此；看，我要将你收归你祖先那里，你必平安地归到他们那里；你的眼睛必看不见我要降在这地方和其中居民身上的一切灾祸。}他因天主的威吓而害怕，哭泣并撕裂衣服，因死亡的催促而免于一切未来的灾祸，因为他将安息在平安中，以致看不见那些事。因此，死者的灵魂所在之处，他们看不见此生中人们所做或所发生的一切。那么，他们如何看见自己的坟墓，或自己的身体，无论它们是丢弃或埋葬？他们如何参与活人的苦难呢？他们要么承受自己的恶果，如果他们有这样的功过；要么安息在平安中，正如应许给这位\Person{约史雅}的，在那里他们不受任何苦，无论是自己受苦，还是与他人同受苦，因他们已从活着时所受的一切苦中解脱出来，无论是自己受苦还是与他人同受。\par

17. 有人也许会说：\Quote{如果死者对生者没有任何关怀，那么那个在地狱里受折磨的富人，如何请求先祖\Person{亚巴郎}派遣\Person{拉匝禄}到他尚未死去的五个兄弟那里，去警告他们，免得他们自己也来到这个受折磨的地方？}但这是否意味着，因为富人说了这话，他就知道他的兄弟们在做什么，或那时在遭受什么？同样，他关心生者，尽管他完全不知道他们在做什么，正如我们关心死者，尽管我们承认不知道他们在做什么。因为如果我们不关心死者，我们就不会像现在这样为他们向天主祈求。最后，\Person{亚巴郎}并没有派遣\Person{拉匝禄}，并且回答说，他们这里有\Person{梅瑟}和先知，他们应当听从他们，免得来到这些折磨中。这里又出现一个问题：此地发生的事，先祖\Person{亚巴郎}自己如何不知道，而他却知道\Person{梅瑟}和先知在这里，也就是说，他们的书，通过顺从这些书，人们可以逃脱地狱的折磨：并且简言之，他知道那个富人活着时享尽了福乐，而穷人\Person{拉匝禄}则生活在劳苦和悲伤中？因为他对他这样说：\BibleQuote{孩子，你应记得你生前享尽了福，而拉匝禄却受尽了苦。}那么他知道这些事，这些事当然发生在活人中间，而不是死人中间。确实如此，但也许，不是在他们活着时发生这些事的时候，而是在他们死后，他通过\Person{拉匝禄}得知了这些事：这样先知所说的\BibleQuote{亚巴郎不认识我们}就不虚假了。\par

因此，我们必须承认，死者确实不知道此间正在发生的事，但在它发生之时；然而其后，他们从那些在此去世后到他们那里的人那里听说；并非所有事，而是那些被允许传述这些事的人（他们也被允许记起这些事）所告知的事；并且那些被告知的人适宜听到这些事。也可能是，从那些存在于此处所发生之事中的天使那里，死者听到一些事情，万物所服从的那一位认为适合让每个死者听到的事情。因为如果没有天使既能在生者之处也能在死者之处，主耶稣就不会说：\BibleQuote{那穷人也死了，并被天使们带到亚巴郎的怀中。}因此，那些从此处将天主愿意的人带到那里的人，能够一时在此、一时在彼。也可能是，死者的灵魂确实学到一些此处发生的事，即必须知道的事，以及必须知道这些事的人；不仅是过去或现在的事，甚至是未来的事，由\OriginalTerm{天主圣神}{Spirit of God}启示给他们：正如并非所有人，而是先知们在此生活时就知道，而且连他们也不是知道一切，而只是天主眷顾认为适宜启示给他们的事。此外，有些死者被派往生者那里，正如另一方面，\Person{保禄}从生者中被提到\Place{乐园}，圣经为此作证。因为\Person{撒慕尔}先知向活着的\Person{撒乌耳}显现，甚至预言了国王将要遭遇的事：虽然有些人认为那并非撒慕尔本人（因他不可能被巫术召来），而是一些适合如此邪恶工作的幽灵冒充了他的模样；不过，\WorkTitle{德训篇}（据说是\Person{息辣之子耶稣}所著，且因其文风相似而被认为是\Person{撒罗满}所作）在赞扬先祖时，称撒慕尔死后仍预言。但如果这本书因\OriginalTerm{希伯来正经}{canon of the Hebrews}（因为不包含在其中）而受到质疑，我们对\Person{梅瑟}又该说什么呢？我们确实读到他在\WorkTitle{申命纪}中去世，而在福音中他与未死的\Person{厄里亚}一同向活着的人显现。\par

由此也解决了那个问题：如果死者不知道生者所做的事，那么殉道者们借着赐予祈祷者的恩惠表明他们关心人间事务，这怎么可能呢？因为不仅借由恩惠的效果，而且就在人们的注视下，众所周知，当蛮族攻击\Place{诺拉}时，宣信者\Person{斐理斯}（你们虔诚地热爱他的本地身份）显现了，我们并非通过不确定的传闻，而是通过可靠的见证人得知此事。然而，这些事情是由天主以远超寻常秩序的方式显示的，而寻常秩序是对每种受造物所分配的。因为水在主愿意时瞬间变成酒，这并不意味我们不应重视水在元素正常秩序中的价值与效能，以区别于那个神圣工作的罕见性，或更确切地说，独特性；同样，拉匝禄复活了，并不意味着每个死人都能随意复活；也不意味着一个活人能像唤醒睡觉的人那样使一个死人复活。人间事物的界限是另一种，神圣德能的标记是另一种：自然发生的事物与奇迹发生的事物是另一种；虽然天主临在于自然界使其存在，而在奇迹中自然界也并非缺席。因此，我们不应认为，仅仅因为有些殉道者临在治病助人，就可以任意关心生者的事务；而应理解为，殉道者关心生者的事务必定是出于神圣的大能，因为凭借其本性，死者是不可能关心生者的事务的。\par

20. 然而，这是一个超出我理解能力的问题：殉道者如何帮助那些确实因他们而得助的人？是他们在同一时间亲自临在于如此不同的地方，且彼此相距遥远，无论是位于他们的纪念处，抑或在纪念处之外，只要感到他们临在的任何地方？或是在他们自身居于与功德相宜的地方，远离与凡人的一切交往，却以普遍的方式为求助者的需要祈祷（如同我们为亡者祈祷，虽然我们不在他们面前，也不知他们在何处或做什么），而无所不在的全能者天主，既不局限在我们内，也不远离我们，俯听并应允殉道者的祈祷，藉着遍及各处天使的服务，赐予人类那些安慰，并且祂认为在此悲惨人生中应当赐予的人，祂就赐予；关于祂的殉道者，祂在何处、何时、以何种方式，尤其藉着他们的纪念处（因为祂知道这对我们有益，为建立对基督的信德，他们因宣认祂而受苦），以奇妙不可言喻的德能和良善，使他们的功德备受尊崇。这是一件我无力达到、过于深奥而无法探究的事；因此，究竟是这两种情形中的哪一种，或者也许两者皆是——有时这些事藉殉道者亲自临在而成，有时藉天使承担殉道者的位格而成就——我不敢断定；我宁愿向知道此事的人求教。因为不能认为无人知晓这些事（当然不是那些自以为知道却不知道的人），因为天主的恩赐有多种，祂赐给这人这个，赐给那人那个，照宗徒所说：\BibleQuote{每人得蒙显示圣神，是为有益于人。这人}，他说，\BibleQuote{藉着圣神获得智慧的言语；另一个人却因同一圣神获得知识的言语；又另一个人在同一圣神内获得信德；另一个人在同一圣神内获得治病的奇恩；有的人能行奇迹，有的人能说预言，有的人能辨别诸灵，有的人能说各种异语，有的人能解释异语。但这一切都是同一圣神所行，随祂的心意个别分给每人。}在宗徒所列举的这些神恩中，谁若获得了辨别诸灵之恩，他就知道这些事，正如它们当被知道的那样。\par

21. 我们可以相信，那位\Person{若望隐修士}就是如此，\Person{老德奥多西皇帝}曾就内战的结果咨询过他；因为他也有预言的恩赐。因为并非每人各自拥有不同的恩赐，而是同一人可拥有多项恩赐，我对此毫无疑义。这位若望，有一次某位极虔诚的妇人想要见他，并通过她的丈夫极力恳求，但他拒绝了这一请求，因为他从未允许妇女见他，却说：\Quote{去，}他说，\Quote{告诉你的妻子，她今夜会见到我，但在她的梦中。}事情果然如此发生：他给了她适合一位已婚信女的劝告。她醒来后，将所见到的那位天主的人（如她丈夫所认识的）以及他对她所说的话，都告诉了丈夫。一位严肃而高贵、最值得信赖的人从他们那里得知此事后又告诉了我。如果我本人见过那位圣洁的隐修士（据说他极有耐心聆听问题，且回答极有智慧），我会就我们的问题向他请教：究竟是他自己来到那妇人的梦中，即他的灵魂以他身体的形状出现，正如我们梦见自己以自己身体的形状出现一样；还是当他自己在做别的事，或者如果睡着了，正在梦见别的事时，要么是通过一位天使，要么以其他方式，那妇人的梦中出现了这样的景象；而他预知此事将如此发生，乃是因预言之神启示了他。因为如果他自己确实在她梦中临现，那么这当然是因着奇迹的恩宠而非本性，是出于天主的恩赐而非他自己的能力。但如果他在做别的事或睡觉而忙于其他景象时，那妇人在梦中看见了他，那么无疑发生了一些类似的事，正如我们在\BibleBook{宗徒}{大事录}中读到的：主耶稣关于\Person{扫禄}向\Person{阿纳尼雅}说话，并告诉他扫禄看见阿纳尼雅前来，而阿纳尼雅自己却不知道。天主的人会根据情况回答我这些问题，然后关于殉道者，我会继续问他：他们是否亲自在梦中或以其他方式临现于那些以他们愿现的形状看见他们的人，尤其是当魔鬼在人们身上承认被殉道者折磨，并请求他们宽恕时？或者这些事是否是通过天使的能力成就的，为光荣和赞美圣人，有益于人们，而圣人们则处于至高安息中，完全自由，从事其他远为优越的景象，远离我们，并为我们祈祷？因为曾在\Place{米兰}，在圣殉道者\Person{伯多削}与\Person{基瓦削}的墓旁，当时在世的\Person{盎博罗削}主教被指名提及，正如已故者被念及他们的名字一样，魔鬼承认他并求他宽恕，而当时他正在忙别的事，对发生的事完全不知。或者这些事有时是通过殉道者的亲临，有时是通过天使的临现而成就的？是否可能区分这两种情况，或通过什么标记区分？或者无人能察觉和判断这些事，除非那人藉天主的圣神拥有那恩赐，\BibleQuote{随祂的心愿，个别分配与人}？我想，这位若望会按我所愿向我讲述这一切：或者通过他的教导我能学到，并知道他所告诉我的乃是真实确定的；或者我会相信我所不知道的，基于他告诉我他所知道的。但如果他可能根据圣经回答，说：\Quote{高于你的事，莫求；比你强的事，莫寻；但主要你做的事，常要思念这些。}我也将感恩地接受。因为如果我们对有些事晦暗不明、不能理解，至少清楚而确定地知道我们不应寻求它们，这并非小的收获；而每个人若想学习某事，认为知道它有益，则应知道不知道它也无害。\par

22. 既然如此，让我们不要认为，为我们所关心的亡者，除了我们隆重恳求的祭台上的祭献、或祈祷、或施舍之外，任何事物能达于他们：虽然并非对为之而行的所有人都得益，而只对那些活着时已赢得使这些事物有益的人。但由于我们无法分辨谁是如此，宜为所有重生者而行，以便没有一个人被遗漏，而这些益处可能且应当达于他们。因为与其缺乏帮助那些需要帮助的人，不如对这些既无妨碍也无帮助的人多行这些事，更胜于此。然而，每人更为自己的近亲朋友殷勤行这些事，以便他们也能同样为他而行。至于埋葬尸体，无论为之耗费多少，并非救恩的帮助，而是人性的职分，基于那情感，即\BibleQuote{没有人曾恨过自己的肉身}。因此，当邻人不再拥有肉身时，他当尽其所能照顾其肉身。若那些不相信肉身复活的人也这样做，那么相信的人更应如此行，以便这种施于身体（虽死但将复活且永远存留）的职分，也在某种程度上成为同一信德的见证。然而，若死者埋葬在殉道者的纪念地，我认为这为亡者如此有益：在将他亦托付于殉道者的转求时，为他恳求的情感得以增加。\par

23. 现在，对于你曾认为适宜询问我的那些事，你得到了我所能给予的答复。如果这答复显得过于冗长，你必须原谅这一点，因为这是出于渴望与你更久交谈的爱。因此，关于这书，你的爱德将如何接纳它，我请求你通过第二封信告知我：尽管它无疑会因它的携带者，即我们的弟兄和同司铎\Person{堪提狄亚努斯}，而更加受欢迎。我已通过你的信认识了他，并全心欢迎了他，且不愿让他离开。因为他在基督的爱德中，以他的临在极大地安慰了我们，而且，说实话，正是因他的催促，我才完成了你的吩咐。因为我的心灵被如此多的事务困扰，若不是他时常提醒我，不让我忘记，那么对你的询问，答复就肯定不会出现了。\par

\SourceRecord{Translated by H. Browne.；Nicene and Post-Nicene Fathers, First Series；Vol. 3.；Edited by Philip Schaff.；Buffalo, NY: Christian Literature Publishing Co.,；1887.；<http://www.newadvent.org/fathers/1316.htm>.}
